\chapter{Графы}

\section{Базовые алгоритмы}

\kactlimport{BellmanFord.h}
\kactlimport{TopoSort.h}
\kactlimport{FloydWarshall.h}


\clearpage
\section{Потоки и разрезы}
	\kactlimport{TeamMinCostFlow.h}
	\kactlimport{TeamDinic.h}
	\kactlimport{MinCut.h}
	\kactlimport{GomoryHu.h}

\section{Паросочетания}

\kactlimport{MinimumVertexCover.h}
\kactlimport{WeightedMatching.h}
\kactlimport{HopcroftKarp.h}
\kactlimport{Kuhn.h}
\kactlimport{GeneralMatching.h}


\section{Алгоритмы на основе DFS}

\kactlimport{SCC.h}
\kactlimport{BiconnectedComponents.h}
\kactlimport{EulerWalk.h}
\kactlimport{2sat.h}


\section{Раскраска}
	\kactlimport{EdgeColoring.h}

\section{Клики и независимые множества}
	\kactlimport{MaximalCliques.h}
	\kactlimport{MaximumClique.h}
	\kactlimport{MaximumIndependentSet.h}

\section{Деревья}
\kactlimport{HLD.h}
\kactlimport{LinkCutTree.h}
\clearpage
\kactlimport{LCA.h}
\kactlimport{BinaryLifting.h}
\kactlimport{CompressTree.h}
\kactlimport{CentroidDecomposition.h}
\kactlimport{DirectedMST.h}

\section{Формулы}
	\subsection{Число остовных деревьев}
		% I.e. matrix-tree theorem.
		% Source: https://en.wikipedia.org/wiki/Kirchhoff%27s_theorem
		% Test: stress-tests/graph/matrix-tree.cpp
		Построить матрицу $N\times N$. Для ребра $a\to b$ выполнить
		\texttt{mat[a][b]--, mat[b][b]++}; для неориентированного также
		\texttt{mat[b][a]--, mat[a][a]++}. Удалить строку и столбец i.
		Определитель равен числу ориентированных остовов с корнем i
		(рёбра направлены от корня); в неориентированном случае корень любой.

	\subsection{Критерий Эрдёша — Галлаи}
		% Source: https://en.wikipedia.org/wiki/Erd%C5%91s%E2%80%93Gallai_theorem
		% Test: stress-tests/graph/erdos-gallai.cpp
		Простой граф со степенями $d_1\ge\dots\ge d_n$ существует тогда и только тогда, когда сумма степеней чётна и для каждого $k=1\dots n$:
		\[ \sum _{i=1}^{k}d_{i}\leq k(k-1)+\sum _{i=k+1}^{n}\min(d_{i},k). \]
